\chapter{Mise en œuvre de la solution}

\section{Architecture cible}

La solution que j'ai conçue s'appuie sur une chaîne d'alerting entièrement pilotée par le code. Les
métriques d'infrastructure sont surveillées via CloudWatch, tandis que les métriques d'application sont
collectées par les collecteurs OpenTelemetry et stockées dans AMP \cite{aws-amp}. Les règles d'alerte,
décrites en Terraform \cite{terraform-docs}, sont déployées automatiquement sous forme d'alarmes
dans CloudWatch et de règles Prometheus \cite{prometheus-docs} dans AMP. Lorsqu'une alerte se
déclenche, elle est transmise à incident.io \cite{incidentio} — depuis AMP via son Alertmanager
managé, depuis CloudWatch via un topic SNS. C'est incident.io qui prend alors en charge le
regroupement (grouping), le routage et l'escalade des alertes, avant de notifier les bonnes personnes
sur le canal le plus adapté (Slack, SMS, appel, notification push ou e-mail). En m'appuyant sur un outil
spécialisé plutôt que sur une chaîne maison, j'ai gagné en fiabilité et en simplicité de maintenance.

\begin{figure}[h]
	\centering
	\fbox{%
		\parbox{0.9\textwidth}{%
			\textbf{Schéma de l'architecture cible (à remplacer par la capture de l'architecture retenue).}\\
			GitHub (Terraform) $\rightarrow$ règles déployées dans CloudWatch (alarmes) et AMP (règles
			Prometheus) $\rightarrow$ transmission à incident.io (via SNS pour CloudWatch, via
			l'Alertmanager managé pour AMP) $\rightarrow$ incident.io : grouping, routage, escalade
			$\rightarrow$ Slack / SMS / Appel / Push / E-mail.
		}%
	}
	\caption{Architecture cible de la chaîne d'alerting}
	\label{fig:archi-alerting-cible}
\end{figure}

\section{Le module Terraform et les fichiers de configuration}

J'ai développé un module Terraform dédié à la gestion centralisée des alertes génériques CloudWatch
et AMP. Les développeurs disposent de deux leviers : modifier les
seuils par défaut en ajustant les valeurs dans le fichier \texttt{.tfvars} spécifique à l'environnement, et
choisir d'appliquer ou non une alerte à leur service. L'arborescence du dossier
\texttt{shared/shared-alarms} attribue à chaque équipe un dossier contenant trois fichiers \texttt{.tfvars}, un par
environnement (production, développement, prelive).

\begin{figure}[h]
	\centering
	\fbox{%
		\begin{minipage}{0.9\textwidth}
			\ttfamily\raggedright
			environment = "dev"\\[0.5em]
			services = \{\\
			\ \ "cloudwatch-alarms-api" = \{\\
			\ \ \ \ cluster\ \ \ \ \ = "clu-infra-staging"\\
			\ \ \ \ cpu\ \ \ \ \ \ \ \ \ = \{ threshold = 78 \}\\
			\ \ \ \ memory\ \ \ \ \ = \{ threshold = 90 \}\\
			\ \ \ \ latency\_p99 = \{ threshold = 2 \}\\
			\ \ \}\\
			\}
		\end{minipage}%
	}
	\caption{Structure d'un fichier \texttt{.tfvars} permettant à une équipe de surcharger les seuils par défaut}
	\label{fig:extrait-tfvars}
\end{figure}

Une fois leurs modifications réalisées, les développeurs soumettent une Pull Request, ce qui déclenche
le processus de CI/CD appliquant les changements. La gestion des alertes personnalisées n'est pas
encore prise en charge par le module ; les développeurs conservent toutefois la possibilité de créer ce
type d'alerte directement dans leur fichier CloudFormation.

\section{La chaîne de routage des notifications}

Une fois qu'une alerte se déclenche dans CloudWatch ou dans AMP, elle est transmise à incident.io, qui
applique une logique de notification complète. J'ai retenu cet outil car il prend en charge, de façon
native et intelligente, les traitements qu'il aurait fallu développer et maintenir soi-même :

\begin{itemize}
	\item \textbf{regrouper} les alertes liées à un même service afin d'éviter la submersion (grouping) ;
	\item \textbf{router} chaque alerte vers les bonnes personnes selon le service impacté, sa criticité et l'environnement ;
	\item \textbf{escalader} automatiquement vers le niveau d'astreinte supérieur toute alerte critique non prise en charge dans les temps ;
	\item \textbf{mettre en sourdine} (silencing) les notifications pendant les maintenances planifiées ;
	\item \textbf{enrichir} chaque notification d'un contexte immédiatement exploitable : niveau de priorité, description de l'alerte, lien vers le runbook et vers le dashboard Grafana du service.
\end{itemize}

Cette étape est décisive : c'est elle qui transforme un signal technique brut en une notification claire,
priorisée et directement actionnable par la personne d'astreinte.

\section{Gestion du cycle de vie des alertes}

Pour rester fiable, le système exige que les alertes créées via Terraform correspondent précisément
aux services réellement déployés : le retrait d'un service doit supprimer automatiquement ses
alertes, et l'ajout d'un service doit les créer automatiquement. Deux approches ont été étudiées :
utiliser EventBridge pour détecter ces évolutions dans les clusters et déclencher une Pull Request de
mise à jour du code Terraform, ou bien se baser sur GitHub comme source de vérité.

À plus long terme, la solution la plus prometteuse consiste à s'appuyer sur le catalogue Backstage
comme source unique de vérité, capable d'identifier les services déployés par chaque équipe et de
cartographier leurs dépendances — par exemple, si le service A dépend de la base de données du
service B, seule une alerte sur B est générée en cas de défaillance de cette base. Le cycle envisagé :
interrogation périodique de l'API Backstage par une GitHub Action, comparaison avec les
fichiers \texttt{.tfvars}, génération automatique d'une Pull Request, puis validation par l'équipe
Infrastructure et application par Terraform.

\section{Stratégie de transition}

La mise en œuvre d'un nouveau modèle d'alerting en production impose une approche incrémentale,
afin de limiter les risques de régression et de faciliter l'appropriation par les équipes. Trois phases ont
été définies :

\begin{itemize}
	\item \textbf{Phase 1 -- Déploiement de la fondation technique} : mise en place de l'infrastructure et du code nécessaires à l'implémentation complète de la solution.
	\item \textbf{Phase 2 -- Coexistence et pilotage (double-run)} : l'ancien système n'est pas immédiatement désactivé ; la nouvelle solution s'exécute en parallèle sur le service d'une équipe pilote, ce qui permet de comparer les alertes émises par les deux systèmes et d'accompagner les développeurs dans l'ajustement des seuils.
	\item \textbf{Phase 3 -- Déclassement de l'ancien système} : une fois la fiabilité de la nouvelle solution démontrée, les alertes de l'ancienne solution sont mises en sourdine, sa visualisation est maintenue temporairement pour observation, avant sa désactivation définitive par la suppression des anciennes Lambdas.
\end{itemize}

Cette trajectoire prudente instaure une confiance progressive dans la nouvelle solution et évite tout
angle mort pendant la bascule.

\section{Ownership et responsabilisation des équipes}

Au-delà de l'aspect technique, la solution corrige l'une des limites majeures du système existant :
l'absence d'ownership des alertes par service. Chaque alerte porte désormais un tag \texttt{Team} et
un tag \texttt{Service} qui l'associent sans ambiguïté à l'équipe responsable, et le routage s'appuie sur
ces tags pour diriger la notification vers le canal Slack de cette équipe plutôt que vers un canal
générique.

Cette responsabilisation passe concrètement par le mécanisme de Pull Request : les équipes ajustent
elles-mêmes, dans leur fichier \texttt{.tfvars}, les seuils et l'activation des alertes de leur service, sous
revue de l'équipe Infrastructure. Chaque alerte critique est en outre associée à un runbook et à un
dashboard Grafana propres au service, ce qui place l'équipe propriétaire en première ligne du
diagnostic — condition nécessaire à une gouvernance durable de l'alerting.
